\section{Offenlegung, Impressum, Nutzungsbedingungen und Hinweise}


\opt{OPTIMIZATION-PRINT}{%
\begin{WidePar}
\FontSansNarrowFamily\small%
\begin{multicols}{3}
}



\minisec{Medieninhaber}

{\noindent\sffamily Arbeitsgemeinschaft Arbeits- und Ausbildungsstandards für
den Sanitätsdienst -- }{\ARGEAASS}
\\
\noindent{}Engerthstraße 146\,/\,8\,/\,13 \\
A-1200 Wien -- AUSTRIA

\begin{compactdesc}
  \item[E-Mail:] office@aass.at
  \item[Homepage:] \url{http://www.aass.at}
  \item[ZVR:] 846678982
\end{compactdesc}

\noindent{}Die {\ARGEAASS} ist ein gemeinnütziger Verein mit Sitz in Wien,
welcher die Entwicklung, Wartung, Qualitätssicherung, Anpassung, Verbreitung und Lehre von
wissenschaftlich und sachlich fundierten Arbeits- und Ausbildungsstandards für
den Sanitätsdienst (AASS), sowie die wissenschaftliche Forschungstätigkeit in
diesem Gebiet bezweckt.

\minisec{Weitere Informationen}

\begin{description}
  \item[Verlags- und Herstellungsort:] Wien
  \item[Publikationsart:] Elektronisch. Die etwaige Herstellung von
  Druckwerken erfolgt durch den autorisierten Nutzer.
  \item[ISBN:] \VarAassIsbn
%   \item[ISSN:] ---
  \item[Status:] {\VarDocVersionTypeName} ({\VarDocVersionTypeDescription})
  \item[Versionsgeschichte der AASS:]~
  
  {\scriptsize %
  \begin{compactdesc}
    \item[0.2.0] \dotfill 2009-07 (Vorabversion)
    \item[0.4.0] \dotfill 2009-09 (Vorabversion)
    \item[0.6.0] \dotfill 2009-10 (Vorabversion)
    \item[0.8.0] \dotfill 2010-01 (Vorabversion)
    \item[0.10.0] \dotfill 2010-06 (Vorabversion)
    \item[0.12.0] \dotfill 2010-08 (Vorabversion)
    \item[0.14.0] \dotfill 2011-01 (Vorabversion)
    \item[0.16.0] \dotfill 2011-03 (Vorabversion)
    \item[0.18.0] \dotfill 2011-08 (Vorabversion)
    \item[0.20] \dotfill 2011-10 (Vorabversion)
    \item[0.22] \dotfill 2012-01 (Vorabversion)
    \item[0.24] \dotfill 2012-08 (Vorabversion)
    \item[0.26] \dotfill 2012-10 (Vorabversion)
    \item[1.0] \dotfill 2013-01 (Release)
    \item[2.0.0] \dotfill 2014-03 (Release)
    \item[2.0.1] \dotfill 2014-08 (Release)
    \item[3.0.0.gamma.1] \dotfill 2014-09 (Vorabversion)
    \item[3.0.0.gamma.2] \dotfill 2015-12 (Vorabversion)
    \item[3.0.0.gamma.3] \dotfill 2016-03 (Vorabversion)
  \end{compactdesc}
  }
  \item[Auszüge:]~
  
  {\scriptsize %
  \begin{compactdesc}
    \item[Megacode] bestehend aus den Kapiteln ASB, ASS, BLS, TOD, IALS
  \end{compactdesc}
  }
  \item[Derivate:]~
  
  {\scriptsize %
  \begin{compactdesc}
    \item[Ernährungspädagogik] 0.2 (2011-11, basierend auf {\AASS} 0.20)
    \item[Ernährungspädagogik] 0.4 (2012-11, basierend auf {\AASS} 0.26)
  \end{compactdesc}
  }
\end{description}

\minisec{Mitarbeit und Meldung von Fehlern}

\noindent{}Wünsche, Anregungen, Beschwerden und Fehlermeldungen nehmen wir
gerne unter 
% \url{http://www.aass.at/bugs}
\href{mailto:anregungen@aass.at}{\nolinkurl{anregungen@aass.at}} entgegen.

\minisec{Nutzungsbestimmungen}

\noindent{} Siehe Abschnitt \enquote*{Lizenzen}.
% Dieses Werk ist nur für den internen Gebrauch gem. §\,42
% Abs.\,6 UrheberG bestimmt. Die Weitergabe ist nicht
% gestattet. 
% 
% Sollten Sie Interesse an den Arbeits- und
% Ausbildungsstandards für den Sanitätsdienst haben, setzen
% Sie sich bitte mit uns in Verbindung.

\minisec{Hinweise}

\noindent{}Wie jede Wissenschaft ist die Medizin ständigen Entwicklungen
unterworfen. Soweit fachliche Angaben gemacht werden, darf der Leser
zwar darauf vertrauen, dass die Autoren große Sorgfalt
darauf verwandt haben, dass diese Angaben dem Wissensstand bei
Fertigstellung des Werkes entspricht. Für Angaben, speziell
über Dosierungsanweisungen und Applikationsformen, kann jedoch
keine Gewähr übernommen werden. 

Jeder Leser ist angehalten, durch
sorgfältige Prüfung der Literatur, der Lehrmeinung so\-wie der
Beipackzettel der verwendeten Präparate und gegebenenfalls nach
Konsultation eines Spezialisten festzustellen, ob die gegebene
Aussage oder Empfehlung für Dosierungen oder die Beachtung von
Kontraindikationen gegenüber der Angabe in diesem
Werk abweicht. Jede Dosierung oder Applikation erfolgt
auf eigene Gefahr des Benutzers. 

Die {\ARGEAASS} bittet jeden Leser, ihm aufgefallene Ungenauigkeiten und
Fehler, sowie Anregungen und konstruktive Kritik, über unsere 
Entwicklungs- und Feedback-Webseite unter
\begin{center}
\url{http://aass.at/entwicklung}
\end{center} 
mitzuteilen.

Geschützte Warennamen und Warenzeichen werden nicht durchgehend besonders
kenntlich gemacht. Aus dem Fehlen eines solchen Hinweises kann also
nicht geschlossen werden, dass es sich um einen freien Warennamen
handelt.

Sofern nicht explizit angegeben, beziehen sich 
geschlechtsspezifische Bezeichnungen -- soweit dies
inhaltlich in Betracht kommt -- auf Frauen und Männer in
gleicher Weise.

\minisec{Die Technik}

\begin{description}
  \item[Satz und Druckvorstufe] \LuaLaTeX{} mit {\sffamily KOMA-Script}
  
  Distribution: \TeX live 2015, mit der AASS
  Style Collection (asc) und vielen, vielen Paketen
  \item[Entwicklungsumgebung] TeXstudio
  \item[Versionsverwaltung] Git
  \item[Literaturverzeichnis] JabRef, BibLaTeX, Biber
\end{description}

\opt{OPTIMIZATION-PRINT}{%
\end{multicols}
\end{WidePar}
}

\bigskip

\minisec{Entwicklungsprozess}

Die Entwicklung der AASS erfolgt in \I[Entwicklungszyklus]{Entwicklungszyklen}, 
sogenannten \T[Sprint]{Sprints}, welche i.\,d.\,R. 3 Monate dauern. Am Anfang 
jedes Sprints werden Entwicklungsziele für den jeweiligen Zyklus vereinbart. Am 
Ende findet ein \I{Review} statt, bei welchem die Entwicklungen des Sprints 
vorgestellt, begutachtet und dieskutiert, sowie anschließend entweder 
angenommen oder verworfen werden. Zuletzt wird ein \I{Änderungsprotokoll} 
erstellt, 
und es findet die Planung für den unmittelbar darauffolgenden Sprint statt. 
Eine Übersicht über die Sprints findet man auf der 
\href{http://aass.at/entwicklung/projects/aass/wiki/Sprints}{E\&F-Homepage} der 
{\AASS}.
Der Entwicklungsprozess der AASS orientiert sich wesentlich an der 
Projektmanagementmethode \href{http://de.wikipedia.org/wiki/Scrum}{Scrum}.


\TableWide
{Versionstypen}%1
{Es gibt sechs verschiedene Versionstypen, 
welche den 
Entwicklungsfortschritt einer Version angeben}%2
{\label{tab:versionstypen}}%3
{}%4
{}%5
{ %\scriptsize
\begin{tabulary}{\linewidth}[H]{llLl}\addlinespace\toprule\addlinespace
\noindent\TabHeading{Kennung}  &
\noindent\TabHeading{Versionstyp} &
\noindent\TabHeading{Beschreibung} &
\noindent\TabHeading{Beispiel}
\\\addlinespace\midrule\addlinespace
\noindent\TabSub{nightly}
&
\E{Laufende Entwicklung}
&
Es handelt sich dabei um eine Zwischenversion ohne formale 
Kriterien zur Fertigstellung.
Die Versionstyp-spezifische Nummerierung ergibt sich aus dem 
Datum.
&
\E{2.0.0.nightly.20140119164900}
\\\addlinespace
\noindent\TabSub{alpha}
&
\I{Zwischenversion}
&
Version ohne wesentliche formale Kriterien zur 
Fertigstellung, 
gedacht als Version zu Demonstrations- oder speziellen Testzwecken. 
Eine Version, 
welche zu einem Review eingereicht wird, 
wäre auch eine Alpha Version.
Die Versionstyp-spezifische Nummerierung ergibt sich aus dem 
Datum.
&
\E{2.0.0.alpha.20140119164900}
\\\addlinespace
\noindent\TabSub{beta}
&
\I{Sprintversion}
&
Version, welche nach einem Sprint 
(Entwicklungszyklus, Iteration) 
fertig gestellt wurde. 
Grundsätzlich wäre sie theoretisch nach formalen Gesichtspunkten 
einsetzbar, 
jedoch entspricht sie nicht den inhaltlichen Qualitätsstandards 
und ist auch nicht zum produktiven Einsatz gedacht.
Die Versionstyp-spezifische Nummerierung ergibt sich aus der 
Nummer der Iteration.
&
\E{2.0.0.beta.44}
\\\addlinespace
\noindent\TabSub{gamma}
&
\I{Vorabversion}
&
Version, 
welche für den produktiven Einsatz in kontrollierten Umgebungen 
(Partnerinstitutionen) gedacht ist.
Die Versionstyp-spezifische Nummerierung ergibt sich aus einer 
fortlaufenden Nummerierung der Vorabversionen.
&
\E{2.0.0.gamma.1}
\\\addlinespace
\noindent\TabSub{delta}
&
\I{Release Candidate}
&
Für die Veröffentlichung 
eingereichte Version.
Die Versionstyp-spezifische Nummerierung ergibt sich aus dem Zusatz 
RC 
und einer durch einen Bindestrich getrennten 
fortlaufenden Nummerierung (RC-1, \ldots).
&
\E{2.0.0.delta.RC-1}
\\\addlinespace
\noindent\TabSub{final}
&
\I{Release}
&
Es handelt sich hierbei um die endgültige Version,
welche für den öffentlichen und produktiven Einsatz konzipiert wurde. 
Die finale Version hat keine Versionstyp-spezifische 
Nummerierung.
&
\E{2.0.0}
\\\addlinespace\bottomrule
\end{tabulary}

\bigskip

Eine \IX{Versionsnummer} besteht aus
\begin{compactenum}
  \item einem dreistelligen \II{Primärschlüssel}, bestehend aus
  \begin{compactenum}
    \item der \II{Hauptversionsnummer} 
    (\I{Major}) 
    zur Kennzeichnung der \II{Hauptversionslinie},
    \item der \II{Nebenversionsnummer} 
    (\I{Minor}) 
    zur Kennzeichnung von Update-Versionen der jeweiligen 
    \I{Hauptversionslinie} und
    \item der \II{Revisionsnummer} für Fehlerbehebung;
  \end{compactenum}
  \item dem \II{Versionstyp} und einer
  \item nachgeordneten 
  \II[Nummerierung!versionstyp-spezifische]{Versionstyp-spezifischen 
  Nummerierung} (\I{VSN}), 
  welche die Eindeutigkeit der Versionskennung sicherstellen soll.
  \item Die volle Versionskennungen beinhaltet noch die 
  \II{Variantenbezeichnung}.
  \item Sollte das Produkt einen Auszug des Gesamtwerkes darstellen, 
  so wird noch eine \II{Auszugsbezeichnung} angefügt.
\end{compactenum}
Die Bestandteile der Versionsbezeichnung werden durch Punkte getrennt.

\begin{Synopsis}
  \SynopsisItem{Schema der Versionsbezeichnung: 
  <\E{Primärschlüssel}>.<\E{Versionstyp}>.<\E{VSN}>.<\E{Variante}>.<\E{Auszug}>}
\end{Synopsis}

Die 3 Stellen werden durch Punkte getrennt.
In der Kurzschreibweise können die Nebenversionsnummer und die 
Revisionsnummer weggelassen werden, sofern die jeweilige und die dazu 
untergeordnete Stelle 0 ist 
(z.\,B. 2.1.0 $\longrightarrow$ 2.1,
2.0.0 $\longrightarrow$ 2; 
2.0.1 wird nicht abgekürzt).
}%6
{}%7

\bigskip

\begin{center}
\sffamily\bfseries\LARGE\color{Base} Wünsche, Anregungen, Kritik?

\Huge\color{DefaultRed}\url{http://aass.at/entwicklung}

\end{center}
